\secaoguia{Introdução}

O laboratório de \nomeunidade\ tem por objetivo complementar e sedimentar o
aprendizado dos conceitos teóricos apresentados em sala de aula. Para tanto, estão
previstas as aulas práticas listadas a seguir:

\begin{enumerate}
	\item Portas Lógicas;
	\item Teoremas Booleanos --- Parte 1;
	\item Circuitos Combinacionais;
\end{enumerate}

No Apêndice~A são apresentadas as normas para uso do laboratório, as quais visam
organizar o uso do mesmo, bem como prover uma maior segurança aos alunos durante a
realização das aulas práticas. Pede-se que seja feita uma leitura cuidadosa dessas
normas.

O material teórico de apoio às práticas é o adotado na disciplina teórica
correspondente \cite{ogata2010,nise2017}. Para os fundamentos de eletrônica digital
empregados nas primeiras aulas, recomenda-se \textcite{tocci2018}; as referências dos
relatórios devem seguir a norma \cite{nbr6023}.

As aulas práticas foram planejadas para serem realizadas no tempo de \qty{1}{\hour} e \qty{40}{\minute}, que
corresponde a duas horas/aula. Exige-se um relatório sobre a prática realizada, que
deve ser entregue na aula subsequente. A organização do relatório, bem como a capa do
mesmo, é apresentada no Apêndice~B.